\section{Joint transition learning (\texorpdfstring{\cref{thm:joint}}{Theorem 1(i)})}
\label{app:joint}

This appendix proves \cref{thm:joint}. The trained parameters are the count readout $\mathbf U$ and gate readout $\mathbf W$
(matrices, Frobenius metric) and the query and key \emph{scalars} $q,k$ (Euclidean metric), along fixed query and key
directions; for these parameters the flow below is the exact gradient flow.
We prove the binary case $K=2$, which is the setting of the main text; \cref{rem:J-generalK} explains what changes for general $K$.
The common symbols are $p$, $\varepsilon$, $\lambda$, and the states, match bits and positions of \cref{def:task,app:construction}.






\subsection{Supplied routing: polynomial transition learning}
\label{app:step-supplied-routing}



\subsubsection{Model and exact features}

We count modulo $K=2$ with a fixed query of probability $p\in(0,1)$ and a uniform start; every non-query letter has the same embedding. Let $\bm x_t=\bm e_{s_{t-1}}$, $\mathbf R\bm e_s=\bm e_{s+1}$, $r_b=\sigmoid\bigl(-(-1)^ba\bigr)$ for $b\in\{0,1\}$ and $\sigmoid(z)=(1+e^{-z})^{-1}$.

\paragraph{Seed.} Fix orthonormal residual directions $\hat{\bm e}_0,\hat{\bm e}_\ell$, a count register $(\hat{\bm e}^{\mathrm c}_s)_{s\in\mathbb Z_K}$, a gate register
$(\hat{\bm e}^{\mathrm g}_s)_{s\in\mathbb Z_K}$, and a unit vector $\hat{\bm e}_{\mathrm h}$ of the head space. The embeddings are
\[
 E(s)=\initscale\,\hat{\bm e}^{\mathrm c}_s,\qquad E(\text{query letter})=\initscale(\hat{\bm e}_0-\hat{\bm e}_\ell),\qquad E(\text{other letters})=\initscale(\hat{\bm e}_0+\hat{\bm e}_\ell).
\]
The gate head is head $1$ of block $2$. At the position $i_t$ that holds $s_{t-1}$ and predicts $s_t$ ($t\ge2$), its fixed mask admits
$i_t$ and the position $j_t$ of the current letter, and its relative bias is zero. Its maps are $\mathbf W_Q=q\,\bm\Pi_Q$, $\mathbf W_K=k\,\bm\Pi_K$ and
$\mathbf W_V=v\,\bm\Pi_V$, with scalars $q,k,v$ and fixed matrices
\[
 \bm\Pi_Q=\textstyle\sum_s\hat{\bm e}_{\mathrm h}(\hat{\bm e}^{\mathrm c}_s)^{\!\top},\qquad \bm\Pi_K=\hat{\bm e}_{\mathrm h}\hat{\bm e}_\ell^{\top},\qquad \bm\Pi_V=\sum_s\hat{\bm e}^{\mathrm g}_s(\hat{\bm e}^{\mathrm c}_s)^{\!\top}.
\]
All other embeddings, heads and MLP weights are zero, and the MLP biases are $\pm\initscale$. The only trainable readout blocks are the
$2\times2$ blocks $\mathbf U$ (count register) and $\mathbf W=w(\mathbf I-\mathbf R)$ (gate register), both initialized at zero.
The query and key train along $\bm\Pi_Q,\bm\Pi_K$; the value map $v\bm\Pi_V$ is fixed (\cref{rem:J-line}).

\paragraph{Training.} The trained parameters are $\bm\vartheta=(\mathbf U,\mathbf W,q,k)$, with $\mathbf W$ on the line $\R(\mathbf I-\mathbf R)$, and
$\ddt{\bm\vartheta}=-\nabla_{\bm\vartheta}\Lstep(\bm\vartheta)$ for the Euclidean (Frobenius) norm on $\R^{2\times 2}\times\R(\mathbf I-\mathbf R)\times\R\times\R$. The value factor $v$
and all other parameters are fixed. The four trained parameters follow this one flow at all times; there are no stages and no stopping
rule. The initialization is
 $\mathbf U(0)=0,\qquad w(0)=0,\qquad \tfrac{1}{2}\initscale\le|q(0)|,\ |k(0)|,\ |v|\le2\initscale,$ 
with arbitrary signs.

\begin{lemma}[Exact features]
\label{lem:J-features}
For every $\bm\vartheta$ and $t\ge2$, the logits of the prediction of $s_t$ are
 $\bm z_t=\initscale \mathbf U\bm x_t+\initscale v\,r_{b_t}\,\mathbf W\bm x_t,\qquad a=\frac{\initscale^2qk}{\sqrt{d_{\mathrm h}}}.$ 
\end{lemma}

\begin{proof}
Block $1$ and all MLPs add zero, so each residual is its embedding. The query at $i_t$, the keys and the values are
\begin{gather*}
 \text{query}_{i_t}=\initscale q\,\hat{\bm e}_{\mathrm h},\qquad \text{key}_{j_t}=(-1)^{b_t}\initscale k\,\hat{\bm e}_{\mathrm h},\qquad \text{key}_{i_t}=0,\\
 \text{value}_{i_t}=\initscale v\,\hat{\bm e}^{\mathrm g}_{s_{t-1}},\qquad \text{value}_{j_t}=0 .
\end{gather*}
The scores are $(-1)^{b_t}a$ at $j_t$ and $0$ at $i_t$, so $i_t$ gets weight $r_{b_t}$. The residual at $i_t$ becomes
 $\initscale\hat{\bm e}^{\mathrm c}_{s_{t-1}}+\initscale v\,r_{b_t}\,\hat{\bm e}^{\mathrm g}_{s_{t-1}},$ 
and the readout gives $\bm z_t$.
\end{proof}

\begin{lemma}[Cyclic symmetry]
\label{lem:J-symmetry}
$\mathbf R\mathbf U(t)=\mathbf U(t)\mathbf R$ for all $t\ge0$; $\mathbf W(t)\in\R(\mathbf I-\mathbf R)$ commutes with $\mathbf R$ by definition.
\end{lemma}

\begin{proof}
Shifting $s_0$ by one preserves the data law and the match bits, and $\mathrm{CE}(\bm z,j+1)=\mathrm{CE}(\mathbf R^{-1}\bm z,j)$ for the cross-entropy. By \cref{lem:J-features},
 $\Lstep\circ S=\Lstep,\qquad S(\mathbf U,\mathbf W,q,k)=(\mathbf R^{-1}\mathbf U\mathbf R,\ \mathbf W,\ q,\ k).$ 
$S$ is a linear isometry, so $S\bm\vartheta(t)$ solves the flow, with $S\bm\vartheta(0)=\bm\vartheta(0)$. By uniqueness, $S\bm\vartheta(t)=\bm\vartheta(t)$.
\end{proof}

\subsubsection{Reduction}

Let $\pi_0=1-p$, $\pi_1=p$, $H_p=-p\log p-(1-p)\log(1-p)$, $p_{\min}=\min(p,1-p)$, $\ell(m)=\log(1+e^{-m})$ and
$\phi(m)=-\ell'(m)=\sigmoid(-m)$.

\begin{lemma}[Reduced loss and flow]
\label{lem:J-reduce}
Write $\mathbf U=u_0\mathbf I+u_1\mathbf R$ and let
 $c=\initscale(u_0-u_1),\qquad g=2\initscale vw,\qquad m_b=(-1)^b\bigl(c+g\,r_b\bigr).$ 
The prediction of $s_t$ with match bit $b_t=b$ has margin $m_b$ for its target over the other count. Hence, for every $n$,
 $\Lstep=\mathcal L(c,g,a):=(1-p)\,\ell(m_0)+p\,\ell(m_1).$ 
With $\qkbalance=q(0)^2-k(0)^2$,
\begin{equation}
 \begin{gathered}
 \ddt{c}=-\initscale^2\pd{\mathcal L}{c},\qquad \ddt{g}=-\initscale^2v^2\pd{\mathcal L}{g},\qquad \ddt{a}=-\frac{\initscale^4S}{d_{\mathrm h}}\pd{\mathcal L}{a},\\
 S:=q^2+k^2=\sqrt{\qkbalance^2+\frac{4d_{\mathrm h}a^2}{\initscale^4}}.
 \end{gathered}
 \label{eq:J-flow}
\end{equation}
\end{lemma}

\begin{proof}
\emph{Logits.} By \cref{lem:J-symmetry}, $\mathbf U$ commutes with the swap $\mathbf R$, so $\mathbf U=u_0\mathbf I+u_1\mathbf R$. For a displayed count $s'$,
$\mathbf W\bm x_t=w(\bm e_{s'}-\bm e_{s'+1})$, and by \cref{lem:J-features} the logits of $s'$ and $s'+1$ are $\initscale u_0+\initscale vr_{b_t}w$ and
$\initscale u_1-\initscale vr_{b_t}w$. Their difference is $c+g\,r_{b_t}$. The target is $s'+b_t$, so its margin is $m_{b_t}$ and the cross-entropy is $\ell(m_{b_t})$.

\emph{Flow.} The flow stays on matrices $u_0\mathbf I+u_1\mathbf R$; since $\|\mathbf U\|_F^2=2(u_0^2+u_1^2)$ and $\|\mathbf W\|_F^2=4w^2$, it acts as
$\ddt{u_j}=-\frac12\pd{\Lstep}{u_j}$ and $\ddt{w}=-\frac14\pd{\Lstep}{w}$. With $\pd{\Lstep}{u_0}=-\pd{\Lstep}{u_1}=\initscale\pd{\mathcal L}{c}$ and
$\pd{\Lstep}{w}=2\initscale v\pd{\mathcal L}{g}$ this gives the equations for $c$ and $g$. For the gate,
\[
 \ddt{q}=-\frac{\initscale^2k}{\sqrt{d_{\mathrm h}}}\pd{\mathcal L}{a},\qquad \ddt{k}=-\frac{\initscale^2q}{\sqrt{d_{\mathrm h}}}\pd{\mathcal L}{a},\qquad \frac{d}{dt}(q^2-k^2)=0.
\]
Differentiating $a=\initscale^2qk/\sqrt{d_{\mathrm h}}$ gives the equation for $a$; $S^2=(q^2-k^2)^2+4q^2k^2$ gives its expression for $S$. The flow exists for
all $t\ge0$: the vector field is smooth, $|\pd{\mathcal L}{c}|,|\pd{\mathcal L}{g}|\le1$, so $c,g$ grow at most linearly, and $|\pd{\mathcal L}{a}|\le\frac{|g|}{4}$,
so $(q,k)$ grows at most exponentially by Gr\"onwall.
\end{proof}

\paragraph{The common form.} In sign form, with $S$ and $C$ the signs of the displayed count and the current letter, the margin
$m_b$ is $C\,(c+g\,r_b)$ and $r_b=\sigmoid(-Ca)=\frac12\bigl(1-C\tanh\frac a2\bigr)$. Hence the step logit is
\[
 S\Bigl(\bigl(c+\tfrac g2\bigr)+\bigl(-\tfrac g2\bigr)\tanh\bigl(\tfrac a2\bigr)\,C\Bigr),
\]
which is \cref{eq:gate-logit} with state readout $c+\frac g2$, gate readout $-\frac g2$ and gate strength $\frac a2=\frac{\initscale^2qk}{2\sqrt{d_{\mathrm h}}}$.

\paragraph{Notation.} $r_0+r_1=1$ and $\tau:=r_1-r_0=\tanh(\frac{a}{2})$. Let
\begin{gather*}
 \mathrm{err}_b=\pi_b\,\phi(m_b),\qquad \mathrm{err}=\mathrm{err}_0+\mathrm{err}_1,\\
 h=\tfrac{1}{2}(m_0-m_1)=c+\tfrac{g}{2},\qquad R=m_0+m_1=-g\tau,\qquad Q=g-\tfrac{v^2}{2}c,
\end{gather*}
the weighted errors, the count-only margin, the rule contrast $R=g(r_0-r_1)$ (it cancels the count-only prior), and a gate--count balance. Let $P_b$ be the
probability that the prediction with match bit $b$ assigns to its target.

\begin{lemma}[Derivatives]
\label{lem:J-geometry}
\begin{enumerate}[label=(\alph*),leftmargin=*,itemsep=2pt,topsep=3pt]
\item $\pd{\mathcal L}{c}=\mathrm{err}_1-\mathrm{err}_0$, $\quad\pd{\mathcal L}{g}=r_1\mathrm{err}_1-r_0\mathrm{err}_0=\tfrac{1}{2}\bigl(\pd{\mathcal L}{c}+\tau\,\mathrm{err}\bigr)$, $\quad\pd{\mathcal L}{a}=g\,\dd{\sigmoid(a)}{a}\,\mathrm{err}$.
\item $0<\phi<1$; $\phi(m_b)=1-P_b$; $m\mapsto\phi(m)$ is decreasing; and $\phi(m-2x)\le e^{2x}\phi(m)$ for $x\ge0$.
\item $\ell(m)\ge-m$.
\item $\ddt{\mathcal L}=-\initscale^2(\pd{\mathcal L}{c})^2-\initscale^2v^2(\pd{\mathcal L}{g})^2-\frac{\initscale^4S}{d_{\mathrm h}}(\pd{\mathcal L}{a})^2\le0$.
\end{enumerate}
\end{lemma}

\begin{proof}
(a) Differentiate the loss and margins in \cref{lem:J-reduce}, using $\ell'=-\phi$ and $r_0+r_1=1$.
(b) $\phi(m)=\sigmoid(-m)=1/(1+e^{m})$ is the probability of the wrong count, and
$\phi(m-2x)/\phi(m)=(1+e^{m})/(1+e^{m-2x})\le e^{2x}$.
(c) $\log(1+e^{-m})\ge-m$. (d) is the chain rule with \cref{eq:J-flow}.
\end{proof}

\begin{lemma}[Loss levels]
\label{lem:J-levels}
Let $\mathcal L^\star=\frac{1}{2}p_{\min}\log2$.
\begin{enumerate}[label=(\alph*),leftmargin=*,itemsep=2pt,topsep=3pt]
\item If $\mathcal L<p_{\min}\log2$, then $P_0,P_1>\frac{1}{2}$: every prediction of $s_t$, $t\ge2$, has positive margin.
\item If $P_0,P_1\ge\frac{1}{2}$, then $\mathrm{err}_0^2+\mathrm{err}_1^2\ge \frac{\mathcal L^2}{32}$.
\item If $\mathcal L>\mathcal L^\star$, then $\mathrm{err}_0^2+\mathrm{err}_1^2\ge \frac{p_{\min}^2}{267}$.
\end{enumerate}
\end{lemma}

\begin{proof}
(a) $\pi_b\ge p_{\min}$, so $-\log P_b=\ell(m_b)<\log2$.

(b) For $P\ge\frac{1}{2}$, $-\log P\le\frac{1-P}{P}\le2(1-P)$. By \cref{lem:J-geometry}(b), $\ell(m_b)\le4\phi(m_b)$, so $\mathcal L\le4\,\mathrm{err}$ and
$\mathrm{err}_0^2+\mathrm{err}_1^2\ge\frac{1}{2}\mathrm{err}^2\ge\frac{1}{32}\mathcal L^2$.

(c) If $P_b<\frac{1}{2}$, then $\mathrm{err}_b\ge\frac{\pi_b}{4}\ge\frac{p_{\min}}{4}$ by \cref{lem:J-geometry}(b); otherwise (b) applies. So
 $\mathrm{err}_0^2+\mathrm{err}_1^2\ge\min\Bigl\{\frac{p_{\min}^2}{16},\ \frac{(\mathcal L^\star)^2}{32}\Bigr\}=\frac{(\log2)^2}{128}\,p_{\min}^2\ge\frac{p_{\min}^2}{267}. \qedhere$ 
\end{proof}

\begin{lemma}[Gradient lower bound]
\label{lem:J-gradient}
$(\pd{\mathcal L}{c})^2+(\pd{\mathcal L}{g})^2\ge\frac{1}{3}\tau^2(\mathrm{err}_0^2+\mathrm{err}_1^2)$. If $v^2\le1$, then
 $\ddt{\mathcal L}\le-\tfrac{1}{3}\initscale^2v^2\tau^2\,(\mathrm{err}_0^2+\mathrm{err}_1^2).$ 
\end{lemma}

\begin{proof}
By \cref{lem:J-geometry}(a), $(\pd{\mathcal L}{c},\pd{\mathcal L}{g})^\top=\mathbf C_{\mathrm{grad}}(\mathrm{err}_0,\mathrm{err}_1)^\top$ with $\mathbf C_{\mathrm{grad}}=\bigl(\begin{smallmatrix}-1&1\\-r_0&r_1\end{smallmatrix}\bigr)$. Its singular values $s_1\ge s_2$
satisfy
 \[
\begin{gathered}
s_1s_2=|\det\mathbf C_{\mathrm{grad}}|=|\tau|,\\ s_1\le\|\mathbf C_{\mathrm{grad}}\|_F=(2+r_0^2+r_1^2)^{1/2}\le\sqrt3,\\ s_2\ge\frac{|\tau|}{\sqrt3} .
\end{gathered}
\] 
The second claim follows from \cref{lem:J-geometry}(d).
\end{proof}

\paragraph{Reflection.} The map $(c,g,a)\mapsto(-c,-g,-a)$ exchanges $m_0$ and $m_1$, and so maps $\mathcal L$ and \cref{eq:J-flow} for $p$ to those
for $1-p$. On parameters it is $(u_0,u_1,w,q,k)\mapsto(u_1,u_0,-w,-q,k)$, i.e.\ $q(0)\mapsto-q(0)$. It preserves $q(0)k(0)(1-2p)$, $\qkbalance$ and the
conclusion of \cref{thm:joint}. Hence we assume $p<\frac{1}{2}$, and set
 $h^\star=\log\frac{1-p}{p}>0,\qquad \mu=\frac{p(1-2p)}{14}.$ 
The initialization is \emph{aligned} if $q(0)k(0)<0$, i.e.\ $a(0)<0$.

\begin{lemma}[Drift identities]
\label{lem:J-drift}
\begin{enumerate}[label=(\alph*),leftmargin=*,itemsep=2pt,topsep=3pt]
\item $\operatorname{sign}\ddt{a}=-\operatorname{sign}g$.
\item $\ddt{Q}=-\frac{1}{2}\initscale^2v^2\,\tau\,\mathrm{err}$, so $Q$ is nondecreasing while $a\le0$.
\item $(1+\frac{v^2}{4})\,g=Q+\frac{v^2}{2}\,h$.
\item $\ddt{h}=-\initscale^2(1+\frac{v^2}{4})\pd{\mathcal L}{c}-\frac{1}{4}\initscale^2v^2\,\tau\,\mathrm{err}$.
\item If $h\le\frac{h^\star}{2}$, then $\pd{\mathcal L}{c}\le0$. If $h\le\frac{h^\star}{4}$ and $|R|\le1$, then $\pd{\mathcal L}{c}\le-\mu$.
\end{enumerate}
\end{lemma}

\begin{proof}
(a) By \cref{eq:J-flow} and \cref{lem:J-geometry}(a), $\ddt{a}=-\frac{\initscale^4S}{d_{\mathrm h}}\,g\,\dd{\sigmoid(a)}{a}\,\mathrm{err}$ with $\dd{\sigmoid(a)}{a},\mathrm{err}>0$ and $S=q^2+k^2>0$ (if $\qkbalance=0$, then $S=\frac{2\sqrt{d_{\mathrm h}}|a|}{\initscale^2}$ and $a$ never reaches $0$ by \cref{rem:J-balance}).

(b), (d) By \cref{eq:J-flow} and \cref{lem:J-geometry}(a), $\ddt{Q}=-\initscale^2v^2\bigl(\pd{\mathcal L}{g}-\frac{1}{2}\pd{\mathcal L}{c}\bigr)$ and
$\ddt{h}=\ddt{c}+\frac{1}{2}\ddt{g}$. (c) Substitute $c=h-\frac{g}{2}$.

(e) $m_1=m_0-2h$ and $\pd{\mathcal L}{c}=p\,\phi(m_1)-(1-p)\,\phi(m_0)$. By \cref{lem:J-geometry}(b),
\begin{alignat*}{2}
 \phi(m_1)&\le e^{2\max(h,0)}\phi(m_0) &\qquad& \why{$\phi$ nonincreasing; $\phi(m-2x)\le e^{2x}\phi(m)$}\\
 \pd{\mathcal L}{c}&\le-\phi(m_0)\bigl((1-p)-p\,e^{2\max(h,0)}\bigr),
\end{alignat*}
and $p\,e^{h^\star}=1-p$ gives the first claim. If $h\le\frac{h^\star}{4}$, then $p\,e^{2\max(h,0)}\le p\,e^{h^\star/2}=\sqrt{p(1-p)}$ and
\begin{alignat*}{2}
 (1-p)-\sqrt{p(1-p)}&=\frac{\sqrt{1-p}\,(1-2p)}{\sqrt{1-p}+\sqrt p}\ge\frac{1-2p}{2} &\qquad& \why{$\sqrt{1-p}\ge\tfrac{1}{\sqrt2}\ge\tfrac{1}{2}(\sqrt{1-p}+\sqrt p)$}\\
 \phi(m_0)&\ge\tfrac{1}{2}\sigmoid\bigl(-\tfrac{h^\star}{4}-\tfrac{1}{2}\bigr) && \why{$m_0=h+\tfrac{R}{2}\le\tfrac{h^\star}{4}+\tfrac{1}{2}$}\\
 &\ge\tfrac{1}{4}e^{-1/2}\Bigl(\frac{p}{1-p}\Bigr)^{1/4}\ge\frac{p}{7} && \why{$\sigmoid(-z)\ge\tfrac{1}{2}e^{-z}$}
\end{alignat*}
using $p^{1/4}\ge p$ and $\frac{1}{4}e^{-1/2}\ge\frac{1}{7}$. Multiply.
\end{proof}

\begin{lemma}[Alignment is kept]
\label{lem:J-align}
Let $t_0\ge0$ with
 $a(t_0)<0,\qquad g(t_0)\ge0,\qquad \bigl(1+\tfrac{v^2}{4}\bigr)g(t_0)\ge\tfrac{v^2}{2}\bigl(h(t_0)-\tfrac{h^\star}{2}\bigr)_+ .$ 
Then for all $t\ge t_0$: $g(t)\ge0$, $a(t)\le a(t_0)$, and $h(t)\ge h_0:=\min\{h(t_0),\frac{h^\star}{2}\}$.
\end{lemma}

\begin{proof}
Let $t'=\sup\{t\ge t_0:a<0\text{ on }[t_0,t]\}>t_0$. On $[t_0,t')$, $\tau<0$ and $\mathrm{err}>0$.

\emph{$h\ge h_0$.} Where $h\le\frac{h^\star}{2}$, $\ddt{h}>0$ by \cref{lem:J-drift}(d), (e). If $h(t_1)<h_0$ for some $t_1\in(t_0,t')$, let
$t_2=\sup\{t\le t_1:h(t)\ge h_0\}$. Then $h<h_0\le\frac{h^\star}{2}$ on $(t_2,t_1]$, so $h(t_1)>h(t_2)=h_0$, a contradiction.

\emph{$g\ge0$.} By \cref{lem:J-drift}(b), (c),
\begin{alignat*}{2}
 \bigl(1+\tfrac{v^2}{4}\bigr)g(t)&=\bigl(1+\tfrac{v^2}{4}\bigr)g(t_0)+\bigl(Q(t)-Q(t_0)\bigr) &\qquad&\\
 &\qquad+\tfrac{v^2}{2}\bigl(h(t)-h(t_0)\bigr)\\
 &\ge\bigl(1+\tfrac{v^2}{4}\bigr)g(t_0)-\tfrac{v^2}{2}\bigl(h(t_0)-h_0\bigr) && \why{$Q\uparrow$, $h\ge h_0$}\\
 &\ge0 && \why{$h(t_0)-h_0=(h(t_0)-\tfrac{h^\star}{2})_+$}
\end{alignat*}

\emph{$a\le a(t_0)$.} By \cref{lem:J-drift}(a), $\ddt{a}\le0$ on $[t_0,t')$. If $t'<\infty$, then $a(t')\le a(t_0)<0$ and, by continuity, $a<0$ on a
neighbourhood of $t'$, contradicting the definition of $t'$. So $t'=\infty$.
\end{proof}

\subsubsection{Conditions and scales}
\label{app:J-constants}

\paragraph{Threshold.} Let
\[
 \initscale^\star=\begin{cases}\dfrac{1}{2}&\text{if }q(0)k(0)<0\ \text{(aligned)},\\[8pt]
 \dfrac{\sqrt{\mathrm{gap}\,p}\;h^\star}{712\,(1+\log 2)}&\text{if }q(0)k(0)>0\ \text{(misaligned)},\end{cases}
\]
and in the misaligned case also $|\qkbalance|\ge\mathrm{gap}\,\initscale^2$. Let $0<\initscale\le\initscale^\star$. In the aligned case only $\initscale\le\frac{1}{2}$ is used.
In the misaligned case, since $\sqrt p\,h^\star\le\sqrt p\log\frac{1}{p}\le\frac{2}{e}$, $p^2(h^\star)^6\le(\frac{3}{e})^6$ and $\frac{p(h^\star)^4}{1-2p}\le10$
(\cref{app:verify}),
\begin{equation}
 \initscale\le\frac{1}{2},\qquad \initscale^4\le\frac{\mu}{16},\qquad \initscale^6\le\frac{p}{16(1+\log 2)},\qquad
 \initscale^2\le\frac{\mathrm{gap}\,p\,(h^\star)^2}{506944}.
 \label{eq:J-eps}
\end{equation}
These are the only smallness properties of $\initscale$ used below.

\paragraph{Scales.} Let
\[
 a_{\mathrm{sc}}=\frac{\initscale^4}{\sqrt{d_{\mathrm h}}},\qquad T_1=\frac{h^\star}{2\initscale^2\mu},\qquad
 T_{\mathrm{cross}}=\frac{15840\sqrt{d_{\mathrm h}}}{\mathrm{gap}\,p\,h^\star\initscale^4},\qquad T_{\max}=\frac{h^\star\sqrt{d_{\mathrm h}}}{16\initscale^6}.
\]
From the seed and $\initscale\le\frac{1}{2}$,
\begin{equation}
 \begin{gathered}
 \tfrac{1}{4}\initscale^2\le v^2\le4\initscale^2\le1,\qquad \tfrac{1}{4}a_{\mathrm{sc}}\le|a(0)|\le4a_{\mathrm{sc}}\le\tfrac{1}{4},\\
 |a|\le1\implies\dd{\sigmoid(a)}{a}\ge\tfrac{1}{6},\qquad 0\le x\le2\implies\tanh(\tfrac{x}{2})\ge\tfrac{x}{4} .
 \end{gathered}
 \label{eq:J-scales}
\end{equation}
In the misaligned case \cref{eq:J-eps} gives $a_{\mathrm{sc}}\le\initscale^4\le\frac{\mu}{16}$, $T_1\le\frac{1}{2}T_{\max}$ (equivalent to $\initscale^4\le\frac{\mu\sqrt{d_{\mathrm h}}}{16}$) and
$T_{\mathrm{cross}}<\frac{1}{2}T_{\max}$ (equivalent to $\initscale^2<\frac{\mathrm{gap}\,p\,(h^\star)^2}{506880}$).

\subsubsection{Phase I: misaligned start}

An aligned start skips this phase. For a misaligned start, we need the count prior to orient the gate before the other parameters
move too far. The crossing lemma follows the count margin and the gate readout together, rather than treating either as fixed.

\begin{lemma}[Crossing]
\label{lem:J-cross}
Let $q(0)k(0)>0$ and $t_{\mathrm a}=\inf\{t\ge0:a(t)\le-\frac{a_{\mathrm{sc}}}{8}\}$. Then $t_{\mathrm a}\le T_1+T_{\mathrm{cross}}$, and
 $g\ge0\ \text{on }[0,t_{\mathrm a}],\qquad |Q(t_{\mathrm a})|\le\frac{v^2h^\star}{16},\qquad a(t_{\mathrm a})=-\frac{a_{\mathrm{sc}}}{8}.$ 
\end{lemma}

\begin{proof}
Let $t'$ be the supremum of $t\le\min\{t_{\mathrm a},T_{\max}\}$ such that $g\ge0$ and $|R|\le1$ on $[0,t]$; $t'>0$, because $t_{\mathrm a}>0$ (as $a(0)>0$), $R$ is continuous with $R(0)=0$, and
$\ddt{g}(0)=\frac{1}{2}\initscale^2v^2\phi(0)\bigl((1-2p)-\tau(0)\bigr)>0$ since $\tau(0)\le2a_{\mathrm{sc}}<1-2p$. Parts (i)--(v) hold on $[0,t')$; (vi) bounds the time.

\smallskip\noindent(i) \emph{Gate strength.} By \cref{lem:J-drift}(a), $a$ is nonincreasing, so $-\frac{a_{\mathrm{sc}}}{8}<a\le a(0)\le4a_{\mathrm{sc}}$:
 $|\tau|\le\tfrac{1}{2}|a|\le2a_{\mathrm{sc}},\qquad \dd{\sigmoid(a)}{a}\ge\tfrac{1}{6}.$ 

\noindent(ii) \emph{Balance.} By \cref{lem:J-drift}(b) and $\mathrm{err}\le1$,
 $\Bigl|\ddt{Q}\Bigr|\le\initscale^2v^2a_{\mathrm{sc}},\qquad |Q(t)|\le\initscale^2v^2a_{\mathrm{sc}}\,t\le\frac{v^2h^\star}{16}\qquad(t\le T_{\max}).$ 

\noindent(iii) \emph{Count margin.} Where $h\le\frac{h^\star}{4}$,
\begin{alignat*}{2}
 \ddt{h}&\ge\initscale^2\mu-\tfrac{1}{4}\initscale^2v^2|\tau| &\qquad& \why{\cref{lem:J-drift}(d), (e)}\\
 &\ge\initscale^2\mu-\tfrac{1}{2}\initscale^2a_{\mathrm{sc}}\ge\tfrac{1}{2}\initscale^2\mu && \why{(i), $v^2\le1$, $a_{\mathrm{sc}}\le\mu$}
\end{alignat*}
Hence $h\ge0$. Let $T_1'$ be the first time in $[0,t')$ at which $h=\frac{h^\star}{4}$, and $T_1'=t'$ if there is none; then $T_1'\le T_1$, and $h\ge\frac{h^\star}{4}$ on
$[T_1',t')$.

\noindent(iv) \emph{Gate contrast.} On $[0,T_1')$, by \cref{lem:J-drift}(c), (ii) and (iii),
 $\frac{d}{dt}\Bigl[\bigl(1+\tfrac{v^2}{4}\bigr)g\Bigr]=\ddt{Q}+\frac{v^2}{2}\ddt{h}\ge\initscale^2v^2\bigl(\tfrac{\mu}{4}-a_{\mathrm{sc}}\bigr)>0,$ 
so $g>0$ on $(0,T_1']$. After $T_1'$,
 $\bigl(1+\tfrac{v^2}{4}\bigr)g\ge\frac{v^2}{2}\cdot\frac{h^\star}{4}-\frac{v^2h^\star}{16}=\frac{v^2h^\star}{16},\qquad g\ge\frac{v^2h^\star}{32}\ge\frac{\initscale^2h^\star}{128} .$ 

\noindent(v) \emph{Rule contrast: $|R|\le\frac{1}{2}$.}
\begin{alignat*}{2}
 p\bigl(h-\tfrac{R}{2}\bigr)&=-p\,m_1\le p\,\ell(m_1)\le\mathcal L(t)\le\mathcal L(0)=\log 2 &\qquad& \why{\cref{lem:J-geometry}(c), (d)}\\
 \bigl(1+\tfrac{v^2}{4}\bigr)g&\le\frac{v^2h^\star}{16}+\frac{v^2}{2}\Bigl(\frac{\log 2}{p}+\frac{1}{2}\Bigr)\le\frac{v^2(1+\log 2)}{p} && \why{(ii), $h^\star\le\tfrac{1}{p}$, $p<\tfrac{1}{2}$}\\
 |R|&=g|\tau|\le\frac{8\initscale^6(1+\log 2)}{p}\le\frac{1}{2} && \why{(i), \cref{eq:J-eps}}
\end{alignat*}

\noindent(vi) \emph{Crossing time.} Both defining conditions are closed, so they hold at $t'$. By (iv), (v) and continuity, $g(t')>0$ (from
$g\ge\frac{v^2h^\star}{32}$ if $t'>T_1'$, and since $(1+\frac{v^2}{4})g$ increases strictly on $[0,t']$ otherwise) and $|R(t')|\le\frac{1}{2}$; so the
conditions extend past $t'$ unless $t'=\min\{t_{\mathrm a},T_{\max}\}$. On $[T_1',t')$, $m_1=-h+\frac{R}{2}\le\frac{1}{4}$, and
\begin{alignat*}{2}
 \mathrm{err}&\ge p\,\phi(m_1)\ge\tfrac{p}{2}\sigmoid(-\tfrac{1}{4})\ge\tfrac{p}{5} &\qquad& \why{\cref{lem:J-geometry}(b)}\\
 -\ddt{a}&=\frac{\initscale^4S}{d_{\mathrm h}}\,g\,\dd{\sigmoid(a)}{a}\,\mathrm{err}\ge\frac{\initscale^6\mathrm{gap}}{d_{\mathrm h}}\cdot\frac{\initscale^2h^\star}{128}\cdot\frac{1}{6}\cdot\frac{p}{5}
 =\frac{\initscale^8\mathrm{gap}\,h^\star p}{3840\,d_{\mathrm h}} && \why{$S\ge|D|\ge\mathrm{gap}\,\initscale^2$, (i), (iv)}\\
 t'-T_1'&\le\frac{33}{8}a_{\mathrm{sc}}\cdot\frac{3840\,d_{\mathrm h}}{\initscale^8\mathrm{gap}\,h^\star p}=T_{\mathrm{cross}} && \why{$a(T_1')\le4a_{\mathrm{sc}}$}
\end{alignat*}
So $t'\le T_1+T_{\mathrm{cross}}<T_{\max}$, hence $t'=t_{\mathrm a}$. The bounds extend to $t_{\mathrm a}$ by continuity.
\end{proof}

\subsubsection{Phase II and the theorem}

The next two lemmas give a persistent readout floor and a time bound for $|a|$ to reach $1$.

Let $T_h=\frac{h^\star}{4\initscale^2\mu}$ and $g_{\mathrm{low}}=\min\{1,\frac{v^2h^\star}{20}\}\ge\frac{\initscale^2}{80}\min\{1,h^\star\}$.

\begin{lemma}[Gate floor]
\label{lem:J-floor}
Let $t_0$ satisfy the hypotheses of \cref{lem:J-align}, with moreover $h(t_0)\ge0$ and $Q(t_0)\ge-\frac{v^2h^\star}{16}$. Then $g(t)\ge g_{\mathrm{low}}$ for all
$t\ge t_0+T_h$.
\end{lemma}

\begin{proof}
On $[t_0,\infty)$, $a<0$, $Q\ge Q(t_0)$ and $h\ge0$ (\cref{lem:J-align}, \cref{lem:J-drift}(b)). At any $s\ge t_0$ with $g(s)<g_{\mathrm{low}}$,
\begin{alignat*}{2}
 \tfrac{v^2}{2}h(s)&=\bigl(1+\tfrac{v^2}{4}\bigr)g(s)-Q(s)<\tfrac{5}{4}\cdot\tfrac{v^2h^\star}{20}+\tfrac{v^2h^\star}{16}=\tfrac{v^2h^\star}{8} &\qquad& \why{\cref{lem:J-drift}(c), $v^2\le1$}\\
 |R(s)|&=g(s)|\tau(s)|<g_{\mathrm{low}}\le1,
\end{alignat*}
so $h(s)<\frac{h^\star}{4}$, and by \cref{lem:J-drift}(d), (e) and $\tau(s)<0$,
 $\ddt{h}(s)\ge\initscale^2\mu,\qquad \frac{d}{dt}\Bigl[\bigl(1+\tfrac{v^2}{4}\bigr)g\Bigr](s)=\ddt{Q}(s)+\tfrac{v^2}{2}\ddt{h}(s)>0 .$ 
Suppose $g(t)<g_{\mathrm{low}}$ for some $t\ge t_0+T_h$, and let $s_0=\sup\{s\in[t_0,t]:g(s)\ge g_{\mathrm{low}}\}$, or $s_0=t_0$ if there is none. On $(s_0,t]$, $g$
increases, so $g(t)>g(s_0)$; if $s_0>t_0$ then $g(s_0)=g_{\mathrm{low}}$, a contradiction. So $s_0=t_0$, and on $(t_0,t]$, $h$ increases at rate at least
$\initscale^2\mu$: $h(t)\ge\initscale^2\mu\,T_h=\frac{h^\star}{4}$, contradicting $h(t)<\frac{h^\star}{4}$.
\end{proof}

\begin{lemma}[Gate growth]
\label{lem:J-grow}
Let $t_0\ge0$ with $a(t_0)\le-\frac{a_{\mathrm{sc}}}{8}$, and suppose $g\ge g_{\mathrm{low}}$ on $[t_0,\infty)$. Then $t_3=\inf\{t\ge t_0:\mathcal L(t)\le\mathcal L^\star\}$ satisfies
 $t_3-t_0\le\frac{4080\sqrt{d_{\mathrm h}}\,\log\frac{8\sqrt{d_{\mathrm h}}}{\initscale^4}}{p_{\min}\min\{1,h^\star\}\,\initscale^4}+\frac{16020\log 2}{p_{\min}^2\initscale^4} .$ 
\end{lemma}

\begin{proof}
Since $g\ge0$, $a$ is nonincreasing (\cref{lem:J-drift}(a)), so $|a|$ is nondecreasing on $[t_0,\infty)$. Let $t_1=\inf\{t\ge t_0:|a(t)|\ge1\}$. On
$[t_0,\min\{t_1,t_3\})$,
\begin{alignat*}{2}
 \mathrm{err}&\ge\bigl(\mathrm{err}_0^2+\mathrm{err}_1^2\bigr)^{1/2}\ge\frac{p_{\min}}{17} &\qquad& \why{\cref{lem:J-levels}(c)}\\
 \frac{d|a|}{dt}&=\frac{\initscale^4S}{d_{\mathrm h}}\,g\,\dd{\sigmoid(a)}{a}\,\mathrm{err}\ge\frac{2\initscale^2|a|}{\sqrt{d_{\mathrm h}}}\cdot\frac{\initscale^2\min\{1,h^\star\}}{80}\cdot\frac{1}{6}\cdot\frac{p_{\min}}{17} && \why{$S\ge\frac{2\sqrt{d_{\mathrm h}}|a|}{\initscale^2}$, $|a|\le1$}\\
 &=\frac{\initscale^4\min\{1,h^\star\}\,p_{\min}}{4080\sqrt{d_{\mathrm h}}}\,|a|,
\end{alignat*}
so $\min\{t_1,t_3\}-t_0\le\frac{4080\sqrt{d_{\mathrm h}}}{\initscale^4\min\{1,h^\star\}p_{\min}}\log\frac{8}{a_{\mathrm{sc}}}$. If $t_1<t_3$, then on $[t_1,t_3)$, $|a|\ge1$ and
\begin{alignat*}{2}
 \ddt{\mathcal L}&\le-\tfrac{1}{3}\initscale^2v^2\tanh^2(\tfrac{1}{2})\,(\mathrm{err}_0^2+\mathrm{err}_1^2) &\qquad& \why{\cref{lem:J-gradient}}\\
 &\le-\frac{\initscale^4p_{\min}^2}{16020} && \why{\cref{lem:J-levels}(c), $\tanh^2\tfrac{1}{2}\ge\tfrac{1}{5}$}
\end{alignat*}
and $t_3-t_1\le\frac{16020\log 2}{p_{\min}^2\initscale^4}$, since $\mathcal L(t_1)\le\log 2$.
\end{proof}

\begin{lemma}[Descent]
\label{lem:J-descent}
Suppose $|a(t)|\ge\frac{a_{\mathrm{sc}}}{8}$ for all $t\ge t_0$, and let $t_3=\inf\{t\ge t_0:\mathcal L(t)\le \mathcal L^\star\}$. Then
\[
 t_3-t_0\le\frac{2.3\cdot10^6\,d_{\mathrm h}}{p_{\min}^2\initscale^{12}},
\]
every prediction of $s_t$, $t\ge2$, has positive margin for all $t\ge t_3$, and
$\mathcal L(t)\le\frac{4\cdot10^5d_{\mathrm h}}{\initscale^{12}(t-t_3)}$ for $t>t_3$.
\end{lemma}

\begin{proof}
For $t\ge t_0$,
\begin{alignat*}{2}
 \tau^2&\ge\tanh^2(\tfrac{a_{\mathrm{sc}}}{16})\ge\frac{a_{\mathrm{sc}}^2}{1024} &\qquad& \why{\cref{eq:J-scales}}\\
 \ddt{\mathcal L}&\le-\tfrac{1}{3}\initscale^2v^2\tau^2\,(\mathrm{err}_0^2+\mathrm{err}_1^2) && \why{\cref{lem:J-gradient}}\\
 &\le-\frac{\initscale^{12}}{12288\,d_{\mathrm h}}\,(\mathrm{err}_0^2+\mathrm{err}_1^2) && \why{$v^2\ge\tfrac{1}{4}\initscale^2$}
\end{alignat*}
While $\mathcal L>\mathcal L^\star$,
\begin{alignat*}{2}
 \ddt{\mathcal L}&\le-\frac{\initscale^{12}p_{\min}^2}{3280896\,d_{\mathrm h}} &\qquad& \why{\cref{lem:J-levels}(c)}\\
 t_3-t_0&\le\frac{3280896\log 2\,d_{\mathrm h}}{p_{\min}^2\initscale^{12}} && \why{$\mathcal L(t_0)\le\mathcal L(0)=\log 2$}
\end{alignat*}
For $t\ge t_3$, $\mathcal L<p_{\min}\log2$, so all margins are positive and $P_0,P_1>\frac{1}{2}$ (\cref{lem:J-levels}(a)), and
\begin{alignat*}{2}
 \ddt{\mathcal L}&\le-\frac{\initscale^{12}}{393216\,d_{\mathrm h}}\,\mathcal L^2 &\qquad& \why{\cref{lem:J-levels}(b)}\\
 \frac{d}{dt}\frac{1}{\mathcal L}&\ge\frac{\initscale^{12}}{393216\,d_{\mathrm h}} && \qedhere
\end{alignat*}
\end{proof}

\begin{proof}[Proof of \cref{thm:joint}]
\emph{Reduction.} By the reflection, $p<\frac{1}{2}$; then $p_{\min}=p$.

\emph{Aligned start.} $a(0)\le-\frac{a_{\mathrm{sc}}}{4}<0$ and $g(0)=h(0)=0$, so \cref{lem:J-align} applies at $t_0=0$: $a(t)\le-\frac{a_{\mathrm{sc}}}{4}$ for all $t$.
\Cref{lem:J-descent} with $t_0=0$ gives positive margins for $t\ge t_3$ and $\Lstep\to0$. \Cref{lem:J-floor} applies at $t_0=0$ ($h(0)=Q(0)=0$), and
\cref{lem:J-grow} with $t_0=T_h$ gives
 $t_3\le T_h+\frac{4080\sqrt{d_{\mathrm h}}\,\log\frac{8\sqrt{d_{\mathrm h}}}{\initscale^4}}{p\min\{1,h^\star\}\,\initscale^4}+\frac{16020\log 2}{p^2\initscale^4} .$ 

\emph{Misaligned start.} By \cref{lem:J-cross}, $a(t_{\mathrm a})<0$, $g(t_{\mathrm a})\ge0$, and
\[
 \bigl(1+\tfrac{v^2}{4}\bigr)g(t_{\mathrm a})=Q(t_{\mathrm a})+\tfrac{v^2}{2}h(t_{\mathrm a})\ge\tfrac{v^2}{2}\bigl(h(t_{\mathrm a})-\tfrac{h^\star}{8}\bigr)\ge\tfrac{v^2}{2}\bigl(h(t_{\mathrm a})-\tfrac{h^\star}{2}\bigr).
\]
With $g(t_{\mathrm a})\ge0$ the left side is at least $\frac{v^2}{2}(h(t_{\mathrm a})-\frac{h^\star}{2})_+$, so \cref{lem:J-align} applies at $t_0=t_{\mathrm a}$ and $a(t)\le-\frac{a_{\mathrm{sc}}}{8}$ for all $t\ge t_{\mathrm a}$. \Cref{lem:J-descent} with $t_0=t_{\mathrm a}$ gives positive margins for $t\ge t_3$ and $\Lstep\to0$. \Cref{lem:J-floor}
applies at $t_0=t_{\mathrm a}$ ($h(t_{\mathrm a})\ge0$ and $Q(t_{\mathrm a})\ge-\frac{v^2h^\star}{16}$ by \cref{lem:J-cross}), and \cref{lem:J-grow} with $t_0=t_{\mathrm a}+T_h$ gives
\[
 t_3\le\underbrace{T_1+T_{\mathrm{cross}}}_{\text{Phase I}}+\underbrace{T_h+\frac{4080\sqrt{d_{\mathrm h}}\,\log\frac{8\sqrt{d_{\mathrm h}}}{\initscale^4}}{p\min\{1,h^\star\}\,\initscale^4}+\frac{16020\log 2}{p^2\initscale^4}}_{\text{Phase II}} .
\]

\emph{The stated bound.} $T_{\mathrm{cross}}\le\frac{15840\sqrt{d_{\mathrm h}}}{\mathrm{gap}\,p\min\{1,h^\star\}\initscale^4}$. Since $h^\star=\log(1+\frac{1-2p}{p})\le\frac{1-2p}{p}$,
 $T_1=2T_h=\frac{7h^\star}{\initscale^2p(1-2p)}\le\frac{7}{p^2\initscale^2}\le\frac{7}{p^2\initscale^4},$ 
and $\log\frac{8\sqrt{d_{\mathrm h}}}{\initscale^4}\le C\log\frac{d_{\mathrm h}}{\initscale}$, which gives the bound of \cref{thm:joint} with $\gamma=\mathrm{gap}$ (misaligned) or $\gamma=1$ (aligned). In the
aligned case there is no Phase I, and $\mathrm{gap}$ does not enter. \Cref{lem:J-descent} also gives $t_3\le T_1+T_{\mathrm{cross}}+\frac{2.3\cdot10^6\,d_{\mathrm h}}{p^2\initscale^{12}}$, a bound
without the factor $\frac{1}{\min\{1,h^\star\}}$.

\emph{Probability.} Let $X=\frac{|q(0)|}{\initscale}$ and $Y=\frac{|k(0)|}{\initscale}$ be independent and uniform on $[\frac{1}{2},2]$. Then $X^2$ has density
$\frac{1}{3\sqrt z}\le\frac{2}{3}$ on $[\frac{1}{4},4]$, so
 $\Pr\bigl(|\qkbalance|<\mathrm{gap}\,\initscale^2\bigr)=\E\,\Pr\bigl(|X^2-Y^2|<\mathrm{gap}\bigm|Y\bigr)\le2\mathrm{gap}\cdot\tfrac{2}{3}=\tfrac{4}{3}\mathrm{gap} .$ 
For either choice of signs and $\initscale$ below the misaligned threshold, the conclusion of \cref{thm:joint} therefore holds with
probability at least $1-\tfrac{4}{3}\mathrm{gap}$. The randomness is over scalar magnitudes, with fixed structured directions.\qedhere
\end{proof}

\paragraph{Why the gap condition.}\label{rem:J-balance} If $|q(0)|=|k(0)|$ and $q(0)k(0)>0$, then $\qkbalance=0$, $S=\frac{2\sqrt{d_{\mathrm h}}|a|}{\initscale^2}$, and
 $\ddt{a}=-f(t)\,a,\qquad f=\frac{2\initscale^2}{\sqrt{d_{\mathrm h}}}\operatorname{sign}(a)\,g\,\dd{\sigmoid(a)}{a}\,\mathrm{err} .$ 
Then $a$ never changes sign, so a misaligned gate stays misaligned and the crossing estimate of \cref{lem:J-cross}(vi) fails. The condition
$|\qkbalance|\ge\mathrm{gap}\,\initscale^2$ excludes this case. An aligned start needs no such condition.

\paragraph{Readout restriction.}\label{rem:J-line}
For $K=2$, restricting the gate readout to the line $\R(\mathbf I-\mathbf R)$ loses only a constant logit, which does not affect the loss.

\paragraph{General $K$.}\label{rem:J-generalK}
For counting modulo $K\ge3$ the same argument goes through with one more coordinate: the count readout acquires a background
level $s$ for the $K-2$ counts that are neither the displayed count nor its successor, the loss becomes
$\log(1+e^{-m}+(K-2)e^{-s-m/2})$, and the metric factors become $2/K$. Every step above has a general-$K$ counterpart, with
$\log K$ in place of $\log2$ and a time bound $\frac{C}{\varepsilon^4}\bigl(\frac{\sqrt{d_{\mathrm h}}\log(d_{\mathrm h}/\varepsilon)}{p_{\min}\min\{1,h^\star\}\gamma}+\frac{K\log K}{p_{\min}^2}\bigr)$;
the complete general-$K$ proof is in the archival version. \Cref{tab:J-ode} integrates the general-$K$ flow for $K=3,4,6$.

\subsection{Scope and matched separation}
\label{app:step-learning-scope}
\label{app:weak-routing-limitations}
\input{sections/thm4_full_statement}

\paragraph{Letter direction.} If $\bm\Pi_K=\hat{\bm e}_{\mathrm h}\hat{\bm e}_0^\top$ instead, both letters have key $\initscale k\hat{\bm e}_{\mathrm h}$ and the gate weight does not depend on $b_t$.
As in \cref{lem:J-reduce}, both letter values then give $R=0$, and $\Lstep\ge \mathcal L_0(h)\ge H_p$ for every $\bm\vartheta$.

\begin{corollary}[Outcome training from the same initialization]
\label{cor:joint-outcome}
Take the architecture and initialization of \cref{thm:joint}, with any fixed mask at the answer position, and $\mixrate$ as in
\cref{eq:lambda}. Gradient flow on $\Lout$, training $(\mathbf U,\mathbf W,q,k)$ or any larger fixed set of blocks, obeys the plateau bound of \cref{thm:outcome-end} for $\mixrate>0$ whenever $\initscale\le c\min(\mixrate,n^{-1/2})$ and the
widths are at most linear in $n$, as at the executor widths of \cref{prop:represent}. This is exactly the condition of \cref{thm:time}.
\end{corollary}

\begin{proof}
After division by $\initscale$, every map and embedding block has operator norm at most $4\sqrt2$; the count embeddings have singular values
$1$, every MLP bias has $|\tilde b_m|=1$, and the readout is zero: this is \cref{ass:init}. A fixed mask keeps every attention row a softmax over
its visible sites (\cref{ass:arch}) and only lowers the numbers of scores and sites. The MLP biases have normalized magnitude $1$, so their norm is $N_b=\sqrt w$, and \cref{rem:end-width} turns
$\initscale\le\initscale_0=c\min(1,1/N_b)$ into $\initscale\le cn^{-1/2}$ for $w=O(n)$. Hence \cref{prop:score-degree}, \cref{prop:network}(i)--(iii),
\cref{lem:score-sum,lem:end-smooth} and \cref{thm:outcome-end} apply unchanged for this fixed projection.
\end{proof}

\begin{remark}[Training family and scope]\label{rem:joint-scope}
The model trains the readouts and the query--key scalars; the value map and the other features are fixed.
\Cref{thm:joint} covers $p\ne1/2$, gives persistent margins and $L\to0$, and admits the matched outcome plateau above.
At executor widths the supplied-routing time is polynomial in $n$, whereas the matched outcome
time is exponential; the separation requires $n$ of order $\log(1/\varepsilon)/\log(1/\lambda)$.
\end{remark}

\paragraph{Numerical check of the reduced flow.} \Cref{tab:J-ode} integrates \cref{eq:J-flow} for aligned and misaligned seeds.
All $162$ runs solve, and solve times scale as $\initscale^{-4.0}$ to $\initscale^{-4.3}$, consistent with the $\initscale^{-4}\log(1/\initscale)$ bound.

\input{sections/j_ode_table}
